\section*{Appendix B: Experiment Set-up \& Extra Results}
\subsection*{Experiment Set-up}
\begin{table}[h]
    \centering
    \caption{Hardware and Software Specifications.}
    \label{tab:system_specs}
    \begin{tabular}{ll}
        \toprule
        \textbf{Component} & \textbf{Specification} \\
        \midrule
        CPU & AMD Ryzen 7 5800X (@ 4.85 GHz) \\
        GPU & NVIDIA GeForce RTX 3080 Ti \\
        RAM & 32 GB \\
        OS & Ubuntu 24.10 (Kernel 6.11.0) \\
        Framework & FL-bench (based on PyTorch) \\
        \bottomrule
    \end{tabular}
\end{table}    
\label{subsec:setup}
\begin{table*}
    \caption{Ablation on the temperature hyperparameter $\tau$ of H-SFP. Results show the F1 Score at $I_c=5$, $I_e=10$. The optimal performance is consistently achieved at $\tau=0.5$.}
    \centering
    \begin{tabular}{l|cc|cc|cc}
        \hline
        \multirow{2}{*}{\textbf{$\tau$}} & 
        \multicolumn{2}{c|}{\textbf{CIFAR-10}} &
        \multicolumn{2}{c|}{\textbf{CIFAR-100}} &
        \multicolumn{2}{c}{\textbf{HAM10000}} \\ \cline{2-7}
        & \textbf{IID} & $dir(0.7)$ & \textbf{IID} & $dir(0.7)$ & \textbf{IID} & $dir(0.7)$ \\
        \hline
        0.2 & 54.28\% & 36.00\% & 48.15\% & 25.10\% & 71.42\% & 30.55\% \\
        \textbf{0.5} & \textbf{65.44\%} & \textbf{45.52\%} & \textbf{55.23\%} & \textbf{36.65\%} & \textbf{78.86\%} & \textbf{38.36\%} \\
        1.0 & 59.77\% & 37.49\% & 52.88\% & 30.12\% & 76.40\% & 34.70\% \\
        1.5 & 60.03\% & 41.99\% & 53.10\% & 32.50\% & 76.85\% & 36.10\% \\
        2.0 & 54.38\% & 40.44\% & 48.65\% & 28.00\% & 72.12\% & 32.45\% \\
        \hline
    \end{tabular}
    \label{tab:temp_ablation}
\end{table*}
\noindent\textbf{Implementation Details.} 
We implemented our proposed framework using the Python programming language and the FL-bench ecosystem\footnote{https://github.com/KarhouTam/FL-bench/tree/master}. The experiments were conducted on a high-performance workstation to ensure reproducibility. The detailed hardware and software specifications are listed in Table~\ref{tab:system_specs}.

\noindent\textbf{Hardware Environment.} 
The simulations were executed on a machine equipped with an AMD Ryzen 7 5800X CPU (8 cores, 16 threads, clocked at 4.85 GHz) and 32 GB of RAM. For model training and aggregation acceleration, we utilized a single NVIDIA GeForce RTX 3080 Ti GPU with 12 GB of VRAM.

\noindent\textbf{Software Environment.} 
The system runs on the Ubuntu 24.10 operating system (Kernel 6.11.0). All deep learning models were implemented using PyTorch within the FL-bench framework.

\subsection*{Ablation on Temperature Hyperparameter}
\label{sec:supp_temp_ablation}

We investigate the contrastive loss temperature $\tau$ in \cref{tab:temp_ablation}. The results show a clear ``Goldilocks'' effect. A low $\tau$ (e.g., 0.2) is too sharp and unstable, while a high $\tau$ (e.g., 2.0) is too soft and provides an insufficient learning signal. We find $\tau=0.5$ consistently achieves the optimal F1 score across all datasets, balancing feature discrimination and training stability.

